\chapter{Introduction and Background Research}
\label{chapter1}

% ──────────────────────────────────────────────────────────────────────────────────────────────────────────
\section{Introduction}

\subsection{Motivation}

Interactive rendering of vegetation is an important element of video games, simulations and architectural
visualisations. The complexity and variation of natural ecosystems poses a significant challenge as no two
plants are the same and real environments contain hundreds or thousands of individual plants. Traditional
approaches require plants to be manually created through the use of 3D modelling software or constructed
using real world data. These solutions are costly, time-consuming, and typically incapable of responding to
environmental parameters such as wind, light conditions or seasonal changes.

\medskip
Procedural generation is widely used in computer graphics for a number of applications including creating
plants, terrain and, cities~\cite{freiknechtSurveyProceduralGeneration2017,zhangSurveyProceduralContent2022}.
The primary methods for procedurally generating plant life are genetic
algorithms~\cite{haubenwallnerShapeGeneticsUsingGenetic2017}, space
colonisation~\cite{runionsModelingTreesSpace,ratulApplicabilitySpaceColonization2019} and Lindenmayer systems
(L-Systems). Genetic algorithms and space colonisation function by generating branches to fill a predefined
space and are best suited to modelling large, mature trees and struggle to model smaller plants.

\medskip
L-Systems were introduced by Aristid Lindenmayer in 1968 to model the development of filamentous
organisms~\cite{lindenmayerMathematicalModelsCellular1968}. They function by iteratively rewriting a string
through the parallel application of a set of production rules. The resulting string can be interpreted using
turtle graphics to construct a 3D model. The iterative nature of L-Systems allows them to generate plants at
varying levels of maturity in a manner which mimics its natural development.

% ──────────────────────────────────────────────────────────────────────────────────────────────────────────
\section{Background Research}

\subsection{Rewriting Systems}

A rewriting system transforms an object by repeatedly applying a set of production rules, substituting each
symbol with a specified replacement. A typical graphical example of a rewriting system is the Koch snowflake
fractal (Figure~\ref{fig:kochflake}). It is constructed by starting with an equilateral triangle and
repeatedly replacing the middle third of each edge with an outward-pointing triangular bump.

\begin{figure}[h]
  \centering
  \def\svgwidth{0.5\textwidth}
  \input{chapters/figures/KockFlake.pdf_tex}
  \caption{Koch snowflake~\cite{commonsSVGVersionKochFlakepng}.}
  \label{fig:kochflake}
\end{figure}

\subsection{Formal Grammars}

A formal grammar, $G$, is a four-tuple $(N, \Sigma, \omega, P)$ where:
\begin{itemize}
  \item $N$ is a finite set of nonterminal symbols
  \item $\Sigma$ is a finite set of terminal symbols
  \item $\omega \in N$ is the start symbol
  \item $P$ is a finite set of production rules
\end{itemize}

Terminal symbols cannot be replaced while nonterminal symbols can be replaced through the application of
production rules. Production rules have the form $V^* N V^* \rightarrow V^*$ where $V = \Sigma \cup N$ and
$V^*$ denotes the set of strings that can be made from the alphabet $V$, meaning that a production maps from
one string containing a nonterminal symbol to another string.

\subsection{L-Systems}

An L-System can be defined as the three-tuple $(V, \omega, P)$ where $V$ is the alphabet containing both
terminal and nonterminal symbols, $\omega$ is the start symbol or axiom, and $P$ is a set of production
rules~\cite{kariSystems1997}. Unlike formal grammars, where at most one rule is applied per derivation step,
L-Systems apply all applicable productions simultaneously to every symbol in the current string. This
parallel rewriting reflects the concurrent nature of biological development, where all cells in an organism
develop at the same time rather than sequentially.

\subsubsection{Deterministic 0L-Systems}

D0L-Systems use deterministic, context-free rules and are the most basic class of L-System. Each symbol has
exactly one production; any symbol not explicitly assigned a rule is given the identity production $a
\rightarrow a$. Starting from the axiom $\omega_0 = \omega$, successive strings $\omega_1, \omega_2, \ldots$
are generated by simultaneously replacing every symbol according to its production. The Fibonacci string can
be constructed by the following D0L-System:
\begin{align*}
  \omega \quad &: A \\
  p_1    \quad &: A \rightarrow AB \\
  p_2    \quad &: B \rightarrow A
\end{align*}
Starting from $A$, each step yields $A \rightarrow AB \rightarrow ABA \rightarrow ABAAB \rightarrow \dots$.

\subsubsection{Stochastic L-Systems}

Stochastic L-Systems allow for productions to have an associated probability. Each production is extended to
a pair $(p_i \rightarrow \chi_i,\, \pi_i)$ where $\pi_i \in [0, 1]$ is its selection probability, with the
constraint that for each symbol $a$ all productions whose left-hand side is $a$ satisfy $\sum_i \pi_i = 1$.
At each derivation step a production is chosen independently for every symbol according to its probability
distribution. The following example models a symbol that terminates or persists with equal probability:
\begin{align*}
  p_1 \quad &: A \xrightarrow{0.5} B \\
  p_2 \quad &: A \xrightarrow{0.5} A
\end{align*}
Stochastic rules introduce natural variation into plant models: different realisations of the same grammar
produce structurally distinct plants of the same species, capturing the individual variation observed in nature.

\subsubsection{Context-Sensitive L-Systems}

Context-sensitive L-Systems introduce productions which are conditionally applied based on the symbols
neighbours. 1L-Systems have productions with one-sided context in the form $a_l < a \rightarrow V^*$ or $a >
a_r \rightarrow V^*$ where $a$ is a nonterminal symbol and $a, a_l, a_r \in V$ while 2L-Systems use
productions in the form $a_l < a > a_r \rightarrow V^*$ and are aware of the symbols to both sides of $a$.
More generally, (k,l)-systems can be described as having a left context of $k$ symbols and a right context of
$l$ symbols. In the case that both a context-sensitive production and context-free production apply to a
symbol the context-sensitive production will take priority. The following example demonstrates rightward
signal propagation:
\begin{align*}
  \omega \quad &: BAAA \\
  p_1    \quad &: B < A \rightarrow B
\end{align*}
Starting from $BAAA$, the signal $B$ propagates one position per step: $BAAA \to BBAA \to BBBA \to BBBB$.
Context-sensitive L-Systems are used in plant modelling to simulate developmental signals, such as flowering
hormones that propagate acropetally through plant tissue, or nutrient signals that affect branching in
adjacent internodes~\cite{prusinkiewiczModelingPlantDevelopment2018}.

\subsubsection{Parametric L-Systems}

Parametric L-Systems attach numerical parameters to symbols, allowing productions to perform arithmetic on
those parameters and pass updated values to successor symbols as well as trigger conditionally depending on
the symbols parameters. Formally, a parametric symbol is written $a(t_1, \dots, t_k)$ where the $t_i$ are
numerical parameters. A parametric production takes the form $a(t_1, \dots, t_k) : \chi(t_1, \dots, t_k) \to
\delta(t_1, \dots, t_k)$ where $\chi$ is a Boolean guard over the parameters and $\delta$ is a string of
parametric symbols whose parameters may be arithmetic expressions over $t_1, \dots, t_k$. A production fires
only when its guard is satisfied. The following example models a countdown symbol:
\begin{align*}
  p_1 \quad &: A(x) : x > 5 \rightarrow B \\
  p_2 \quad &: A(x) : x \leq 5 \rightarrow A(x+1)
\end{align*}
$A$ increments its parameter each step until the counter exceeds five, at which point it terminates as $B$.
Parametric rules are used to model internode elongation (a segment grows each step up to a maximum length),
branching delays (a node waits a fixed number of steps before producing lateral branches), and
resource-driven growth variation.

% ──────────────────────────────────────────────────────────────────────────────────────────────────────────
\subsection{Turtle Graphics}

Turtle graphics provide a method for graphically interpreting the string produced by an L-System. In 2D the
turtle consists of a position and heading and is controlled with commands relative to its current position.
To extend this to 3D the turtle additionally stores a normal which defines the plane the heading is rotated in.

\medskip
The basic instructions are~\cite{prusinkiewiczAlgorithmicBeautyPlants1990, mishraLSystemFractals2007}:
\begin{itemize}
  \item \texttt{F}: move forward, producing a branch segment
  \item \texttt{+}/\texttt{-}: turn left/right by $\delta$ around the normal
  \item \texttt{\&}/\texttt{\^{}}: pitch up/down by $\delta$ around the binormal (normal $\times$ heading)
  \item \texttt{/}/\texttt{\textbackslash}: roll clockwise/counterclockwise by $\delta$ around the heading
  \item \texttt{[}/\texttt{]}: push/pop the turtles state (position, heading, normal)
\end{itemize}

Using turtle graphics with $\delta = 60^\circ$, the Koch snowflake (Figure~\ref{fig:kochflake}) can be
produced by the following L-System:
\begin{align*}
  \omega \quad &: F{-}{-}F{-}{-}F \\
  p_1    \quad &: F \rightarrow F{+}F{-}{-}F{+}F
\end{align*}
The axiom $F{-}{-}F{-}{-}F$ traces an equilateral triangle; after one derivation step, each edge is replaced
by a bump, producing a six-pointed star.

\subsection{Plant Modelling}

L-Systems provide a framework for modelling the development of branching plants. A grammar encodes branching
topology, internode geometry, and leaf production as production rules. Iterative application generates the
successive developmental stages of the organism. The resulting symbol string is then interpreted by a turtle
to produce a 3D geometric representation.

\medskip
The choice of grammar class determines the expressive power of the model. Deterministic context-free grammars
produce regular, self-similar structures. Stochastic grammars introduce per-branch variation by sampling
production probabilities at each rewriting step, yielding the natural irregularity of real plants. Parametric
grammars encode continuous quantities such as internode length, branch angle, and developmental age as symbol
parameters, enabling growth curves to be modelled with a few rules. Context-sensitive grammars allow a symbol
to condition its rewriting on the identity of adjacent symbols in the string, capturing inter-organ signals
such as the acropetal propagation of flowering in an inflorescence~\cite{prusinkiewiczAlgorithmicBeautyPlants1990}.

\subsubsection{Environmental Response}

Plant growth is influenced by a wide range of environmental factors including: gravity (gravitropism), light
direction (phototropism), mechanical stimuli (thigmotropism), wind loading, water availability, and seasonal
temperature cycles. Accurate simulation of all these mechanisms for a specific species requires detailed
parameterisation from field data and is primarily the domain of botanical modelling. For many real-time
applications, such as video games and simulations, fully capturing these complex biological processes is
unnecessary. The primary focus in these domains is for vegetation to be responsive to observable
environmental changes in an intuitive way; this prioritises features such as seasonal response and tropism
over biological accuracy.

\medskip
Tropism is the most widely modelled environmental response in L-System literature. Typically, tropism is
simulated by applying a correction to the turtles heading at each branch segment, bending the growing tip
towards or away from an environmental factor~\cite{prusinkiewiczAlgorithmicBeautyPlants1990}. More complex
models account for proprioception, the sense of positioning, and incorporate species specific
weights~\cite{bastienUnifiedModelShoot2015}.

\medskip
Seasonal variation is typically implemented by scaling rule parameters with time-of-year functions or by
gating productions on the current season. Mech and Prusinkiewicz~\cite{mechVisualModelsPlants1996}
demonstrated this by encoding dormancy as a reduction in internode elongation rate and leaf production in
autumn. A deciduous tree, for example, can be modelled with a single grammar whose shedding probability ramps
up through autumn and recovers in spring, producing realistic bare-winter and leafy-summer silhouettes from a
single set of rules. Phenological events such as flowering or fruiting can similarly be gated by seasonal
parameters computed within parametric productions.

\subsubsection{Positioning}

To create visually convincing scenes the positioning and distribution of plants must be considered in
addition to their form. The accurate, realistic placements of plants within an ecosystem is a complex problem
which requires the simulation of environmental conditions as well as interactions between plants and between
plants and their environment~\cite{chngRealisticPlacementPlants2011}. A comprehensive model for simulating
ecosystem development is beyond the scope of this project. Instead, a simplified kernel-based global-to-local
model is adopted, following the approach of Deussen et al.~\cite{deussenRealisticModelingRendering1998} and
Lane and Prusinkiewicz~\cite{laneGeneratingSpatialDistributions2002}, which captures the key spatial
attributes of plant communities, clustered or overdispersed distributions, species dominance, and edge
effects, without requiring a full ecological simulation.

\medskip
Lane and Prusinkiewicz describe two classes of models for spatial distribution: local-to-global models, which
focus on individual interactions between plants, and global-to-local models, which utilise a distribution of
plant densities to determine individual positions~\cite{laneGeneratingSpatialDistributions2002}. The primary
features of local-to-global models are self-thinning and succession. Self-thinning occurs between similar
plants and describes the competition between plants for resources, with smaller or weaker plants being
outcompeted by larger plants. Succession models changes over time with species being replaced by those with a
higher shade tolerance as an ecosystem matures. Lane and Prusinkiewicz's global-to-local model uses a
probability density function (PDF) over the scene plane which is deformed by kernel functions centred on each
placed plant. An inhibitory kernel reduces the probability of placing another plant nearby whilst a
promotional kernel increases it allowing for the simulation of clustering. Sampling from the cumulative PDF
via inverse transform sampling then produces distributions that match ecologically observed spatial
patterns~\cite{laneGeneratingSpatialDistributions2002}.

% ──────────────────────────────────────────────────────────────────────────────────────────────────────────
\section{Project Aim, Objectives, and Deliverables}

\subsection{Project Aim}

The aim of this project is to implement a procedural vegetation system using L-Systems that generates
realistic and varied plant structures in real time, with support for environmental response and interactive
parameter adjustment.

\subsection{Objectives}

\begin{enumerate}
  \item \textbf{Implement a robust L-System interpreter} utilising turtle graphics to produce geometric
    structures, supporting deterministic, stochastic, parametric and context-sensitive grammar types.
  \item \textbf{Develop L-System grammars for at least three distinct plant types}, with grammars chosen to
    demonstrate different grammar classes and botanical features.
  \item \textbf{Integrate environmental response mechanisms} into the modelling system, allowing plant
    structures to respond to simulated light direction, gravity, wind, and seasonal growth cycles.
  \item \textbf{Implement an interactive 3D visualisation} supporting real-time parameter adjustment and 3D
    mesh generation with level-of-detail management.
\end{enumerate}

\subsection{Deliverables}

The deliverables of this project are:
\begin{enumerate}
  \item A real-time procedural vegetation system capable of rendering scenes of hundreds of individual plants
    at interactive frame rates.
  \item An L-System interpreter supporting deterministic, stochastic, parametric, context-sensitive, and
    combined grammar classes, together with a turtle graphics engine producing 3D branch and leaf geometry.
  \item A collection of L-System grammars for several plant types with accompanying documentation detailing
    parameter descriptions and derivation sequences displaying growth sequences.
  \item An interactive GUI supporting real-time adjustment of all plant, scene, and environment parameters,
    with immediate visual feedback.
  \item A software repository containing the source code.
\end{enumerate}
